\documentclass[11pt]{article}

\usepackage{amsmath,amssymb,amsthm,mathtools}
\usepackage{url}
\usepackage[margin=1in]{geometry}

\newtheorem{theorem}{Theorem}[section]
\newtheorem{lemma}[theorem]{Lemma}
\newtheorem{proposition}[theorem]{Proposition}
\newtheorem{corollary}[theorem]{Corollary}
\newtheorem{remark}[theorem]{Remark}

\newcommand{\Gap}{\Gamma_H}

\begin{document}

\section{Purpose}

For a finite graph $H$, put
\[
 \Gamma_H(W)=t(H,W)-\Phi_H(p(W)),
 \qquad
 \Phi_H(p)=(1-p)^{v(H)}
 \chi_H\!\left(\frac1{1-p}\right).
\]
This note proves a full arbitrary-fibre second-order result on five flat
threshold faces at the balanced bipartite graphon $T_2$, proves the complete
residual axis-profile interaction theorem and then a uniform
arbitrary-profile small-support theorem for four further critical rows, and
resolves the sixth fully flat row, Atlas $152$, by an exact fractional-fibre
counterexample.  It strengthens the
coarse-profile calculation in
\path{notes/turan_local_and_high_density_negative_tests.tex}:
the connections from the exceptional set to either Turan part may now be
arbitrary measurable functions, rather than constants on the bulk parts.

The positive conclusions are local: the five flat faces are controlled at
second order, while one-vertex and rooted-moment positive subintegrals give
a uniform small-support radius for every measurable profile on the four
critical rows.  They prove neither positivity of the graphs nor a
cut-distance neighbourhood theorem.  The
Atlas $152$ conclusion is stronger, because the negative quadratic direction
is realized by one fixed rational five-step graphon and hence gives a full
counterexample rather than merely an infinitesimal obstruction.

\section{The five graphs and their targets}

All graphs have vertex set $\{0,1,2,3,4,5\}$ and chromatic number three.
Their identifying data are as follows.
\[
\begin{array}{c|c|c|l}
\text{Atlas}&\text{graph6}&e(H)&E(H)\\ \hline
130&\texttt{E\char`\{CW}&7&01,02,03,12,34,35,45\\
153&\texttt{Ehf\_}&8&01,04,05,12,15,23,25,34\\
171&\texttt{Ehew}&9&01,04,05,12,23,25,34,35,45\\
174&\texttt{EtTg}&9&01,02,03,14,15,23,25,34,45\\
188&\texttt{EzNG}&10&01,02,05,12,13,15,23,24,34,45
\end{array}
\]
Deletion--contraction, or the exact finite colouring enumeration in the
verifier, gives
\[
\begin{array}{c|l|l}
H&\chi_H(x)&\Phi_H(p)\\ \hline
130&x(x-1)^3(x-2)^2&p^3(2p-1)^2\\
153&x(x-1)(x-2)(x^3-5x^2+9x-7)&
 p(2p-1)(7p^3-12p^2+8p-2)\\
171&x(x-1)(x-2)(x^3-6x^2+13x-11)&
 p(2p-1)(11p^3-20p^2+13p-3)\\
174&x(x-1)(x-2)(x^3-6x^2+14x-13)&
 p(2p-1)(13p^3-25p^2+17p-4)\\
188&x(x-1)(x-2)(x^3-7x^2+18x-17)&
 p(2p-1)(17p^3-33p^2+22p-5).
\end{array}
\]
In particular $\Phi_H(1/2)=0$.  For the last four rows,
\begin{equation}
 \Phi_H'(1/2)=-\frac18.
 \label{eq:negative-slope}
\end{equation}

\section{Arbitrary-fibre perturbations}

Let $A,B,S$ be probability spaces.  For small $\varepsilon>0$, give their
copies masses $1/2-\varepsilon$, $1/2$, and $\varepsilon$, respectively.
Keep the value one on $A\times B$ and zero on $A^2$.  Let
\[
 F:S\times A\longrightarrow[0,1],\qquad
 G:S\times B\longrightarrow[0,1],\qquad
 R:S^2\longrightarrow[0,1]
\]
be measurable, with $R$ symmetric.  These are the exceptional-to-bulk and
exceptional-internal parts of the graphon.  Define, for almost every $z\in S$,
\[
 u(z)=\int_A F(z,a)\,da,
 \qquad
 v(z)=\int_B G(z,b)\,db.
\]
For a pair $z,z'\in S$, the two shared-neighbour moments are
\[
 h_0(z,z')=\int_A F(z,a)F(z',a)\,da,
 \qquad
 h_1(z,z')=\int_B G(z,b)G(z',b)\,db.
\]
The sharp Frechet bounds are
\begin{equation}
 \max(0,u+u'-1)\leq h_0\leq\min(u,u'),
 \qquad
 \max(0,v+v'-1)\leq h_1\leq\min(v,v').
 \label{eq:frechet}
\end{equation}
Indeed, the lower bound follows by integrating
$(1-F(z,\cdot))(1-F(z',\cdot))\geq0$, and the upper bound follows from
$FF'\leq F,F'$.  The proof for $G$ is identical.

For Atlas 130 the flat admissible face is
\[
 \Delta_+=\{(u,v)\in[0,1]^2:u+v\geq1\}.
\]
For Atlas 153, 171, 174, and 188 the repaired flat face is
\[
 \Delta_-=\{(u,v)\in[0,1]^2:u+v\leq1\}.
\]
These are exactly the zero faces of the supporting first variation.  In the
normalization of the exact audit, the density direction is
$D(u,v)=u+v-1$.  The direct supporting polynomial is identically zero for
Atlas 130.  For the other four rows it is $8D$, while the repaired supporting
polynomial is identically zero.  Thus the displayed triangles exhaust the
flat branches; no profile has been selected within them.

On the four repaired rows, put a single constant weight $y_\varepsilon$ on
$B^2$, chosen from the aggregate edge density so that $p=1/2$.  When the
average first-order density displacement is negative,
$y_\varepsilon\geq0$ and $y_\varepsilon=O(\varepsilon)$, hence this is a
valid graphon for all sufficiently small $\varepsilon$.  If that average is
zero and the uncorrected second-order displacement is positive, no repair is
needed: by \eqref{eq:negative-slope}, $\Phi_H(p)<0$ just above $1/2$, so
$t(H,W)\geq0$ already excludes a negative direction.  Atlas 130 has zero
diagonal first variation, so an $O(\varepsilon^2)$ density repair on its
boundary does not alter the quadratic coefficient.

\section{The exact pair kernel}

We give a finite definition that is convenient both for review and for
formalization.  Temporarily replace two infinitesimal portions of $S$ by
classes $L,R$ of masses $s,t$.  The bulk masses are
\[
 m_A=\frac12-s-t,\qquad m_B=\frac12.
\]
The macro-cell values are
\[
 W_{AB}=1,\quad W_{AA}=0,\quad
 W_{LA}=u,\quad W_{LB}=v,\quad
 W_{RA}=u',\quad W_{RB}=v',\quad W_{LR}=r.
\]
For a bulk vertex adjacent to both exceptional classes, replace the formal
product $uu'$ by $h_0$ on $A$ and $vv'$ by $h_1$ on $B$.  This is exactly
the integral over that shared bulk variable.  On a repaired row, set
\[
 W_{BB}=4\left(\frac12-p_{\rm raw}(s,t)\right),
\]
which is the one common repair for the combined two-class perturbation.

For each map $c:V(H)\to\{A,B,L,R\}$ using each of $L,R$ at most once,
multiply its class masses and its edge-cell factors, using the preceding
shared-neighbour rule, and sum.  Let $P_H(s,t)$ be the resulting exact
polynomial after subtracting the chromatic target.  Define
\begin{equation}
 K_{H,r}(u,v,u',v',h_0,h_1)
 :=\frac12[st]P_H(s,t).
 \label{eq:pair-kernel}
\end{equation}
The factor $1/2$ is the usual polarization factor.  If $z=z'$ and
$h_0=\int F(z,a)^2da$, $h_1=\int G(z,b)^2db$, then
$K_{H,r}(z,z)$ is exactly the one-profile quadratic coefficient with
self-cell value $r$.

Because a term of order $\varepsilon^2$ uses at most two exceptional
vertices, Fubini's theorem and polarization give the full coefficient
\begin{equation}
 [\varepsilon^2]\Gamma_H(W_\varepsilon)
 =\int_{S^2}K_{H,R(z,z')}
 \bigl(u(z),v(z),u(z'),v(z'),h_0(z,z'),h_1(z,z')\bigr)
 \,dz\,dz'.
 \label{eq:integrated-kernel}
\end{equation}
All functions are bounded, so every integral is absolutely integrable and
there is no limiting or denominator issue.  Formula \eqref{eq:pair-kernel}
is reconstructed by exact rational arithmetic in
\path{codes/needle_higher_order_tests.py}; the implementation multiplies
only the four bidegrees $1,s,t,st$, so discarded higher powers cannot affect
the result.

\section{Frechet triangulation and certificates}

The coefficient of every positive-degree monomial in $h_0,h_1$ in each of
the ten endpoint kernels ($r=0,1$ for five graphs) is a polynomial in
$u,v,u',v'$ with nonnegative coefficients.  Hence every kernel is
coordinatewise nondecreasing in the two overlaps.  By
\eqref{eq:frechet}, it is enough to substitute their sharp lower endpoints.

Both $\Delta_-^2$ and, after the deficit substitution
\[
 a=1-v,\quad b=1-u,\quad c=1-v',\quad d=1-u',
\]
$\Delta_+^2$ become
\[
 \mathcal D=\{a,b,c,d\geq0:a+b\leq1,\ c+d\leq1\}.
\]
The two Frechet switches cut $\mathcal D$ into six four-simplices.  Put
\[
\begin{array}{lll}
q_0=(0,0,0,0),&q_1=(0,0,0,1),&q_2=(0,0,1,0),\\
q_3=(0,1,0,0),&q_4=(0,1,1,0),&q_5=(1,0,0,0),\\
q_6=(1,0,0,1).
\end{array}
\]
The central region $a+c\leq1$, $b+d\leq1$ is the union of
\[
 (q_1,q_2,q_4,q_5,q_6),\quad
 (q_1,q_3,q_4,q_5,q_6),\quad
 (q_1,q_2,q_3,q_4,q_5),\quad
 (q_0,q_1,q_2,q_3,q_5).
\]
The two side regions are the simplices
\[
\begin{split}
 &(0,0,1,0),(0,1,1,0),(1,0,0,0),(1,0,0,1),(1,0,1,0),\\
 &(0,0,0,1),(1,0,0,1),(0,1,0,0),(0,1,1,0),(0,1,0,1).
\end{split}
\]
This list really is a partition.  For example, the barycentric coordinates
of the four central simplices, in the vertex orders displayed above, are
\[
\begin{array}{l}
(1-a-c,c-b,b,1-c-d,a+c+d-1),\\
(1-a-b,b-c,c,1-b-d,a+b+d-1),\\
(d,1-a-b-d,1-a-c-d,a+b+c+d-1,a),\\
(1-a-b-c-d,d,c,b,a).
\end{array}
\]
If $a+b+c+d\leq1$, the fourth row applies.  Otherwise, either both middle
partial sums are at most one and the third row applies, or comparison of
$b,c$ selects the first or second row.  The side barycentric coordinates are
\[
 (1-a-b,b,1-c-d,d,a+c-1)
\]
and its simultaneous $a\leftrightarrow b$, $c\leftrightarrow d$ image.
Thus all coordinates are nonnegative precisely on the asserted regions.

On each simplex substitute its five barycentric coordinates and homogenize
to the total degree.  The following table records, for the six simplices in
the order above, the total degree, number of nonzero coefficients, and least
positive coefficient.  Every omitted homogeneous coefficient is zero; there
are no negative coefficients and no degree elevation is used.
\[
\begin{array}{c|c|c|c|c}
H&r&\deg&(\#_1,\ldots,\#_6)&(m_1,\ldots,m_6)\\ \hline
130&0&4&(51,51,57,66,37,37)&(1/4,1/4,1/4,1/2,1/4,1/4)\\
130&1&4&(51,51,57,66,37,37)&(1/2,1/2,1/2,3/4,1/2,1/2)\\
153&0&4&(50,50,57,66,36,36)&(1/4,1/4,1/4,1/4,1/8,1/8)\\
153&1&4&(50,50,57,66,36,36)&(1/4,1/4,1/4,1/4,1/4,1/4)\\
171&0&5&(49,53,39,24,24,43)&(3/8,3/8,3/8,3/8,3/8,3/8)\\
171&1&5&(49,53,39,24,37,49)&(3/8,3/8,3/8,3/8,1/8,1/8)\\
174&0&4&(14,15,10,6,6,12)&(3/4,3/4,3/4,3/4,3/4,3/4)\\
174&1&4&(19,21,14,7,13,19)&(3/8,3/8,3/8,3/8,3/8,3/8)\\
188&0&5&(49,53,39,24,24,43)&(1/8,1/8,1/8,1/8,1/8,1/8)\\
188&1&5&(49,53,39,24,37,49)&(1/8,1/8,1/8,1/8,1/16,1/16)
\end{array}
\]
The four central simplex volumes sum to $1/6$; each side simplex has volume
$1/24$.  Their total $1/4$ is the volume of the product of two unit
triangles, providing an additional exact partition check.

\begin{theorem}[Five arbitrary-fibre flat threshold cones]
For $H$ equal to Atlas $130,153,171,174$, or $188$, let
$W_\varepsilon$ be an arbitrary-fibre perturbation described above, lying
almost everywhere in the corresponding flat profile triangle.  Whenever the
raw density is below $1/2$, use the one common aggregate density repair; on
the Atlas $130$ boundary this does not change the quadratic coefficient.
Then
\[
 [\varepsilon^2]\Gamma_H(W_\varepsilon)\geq0.
\]
More strongly, the integrand in \eqref{eq:integrated-kernel} is nonnegative
for almost every pair of exceptional points in the formal repaired
expansion.  On the last four rows, if a boundary perturbation instead has raw
density above $1/2$, it already satisfies $\Gamma_H(W_\varepsilon)\geq0$ for
all sufficiently small $\varepsilon$ because $\Phi_H$ is then negative.
\end{theorem}

\begin{proof}
The endpoint certificates prove $K_{H,0}\geq0$ and $K_{H,1}\geq0$ after
the worst feasible overlap substitution on every one of the six Frechet
regions.  Monotonicity in $h_0,h_1$ extends this to every feasible overlap.
The kernel is affine in $r$, so it is nonnegative for every
$r\in[0,1]$.  Apply this pointwise with $r=R(z,z')$ and integrate
\eqref{eq:integrated-kernel}.
\end{proof}

\section{Four critical rows with an axis zero set}

Atlas $137,139,145$, and $148$ are not fully flat.  Instead, their complete
one-microclass threshold variation has a simple exact zero set.  In the
preceding exceptional-class model, let one class of mass $\varepsilon$ have connection
means $(u,v)\in[0,1]^2$, and use the least constant diagonal repair whenever
$u+v<1$.  With the $2^{v(H)}$ normalization of the exact needle audit, both
the direct and repaired first variations are
\begin{equation}
 2^{v(H)}[\varepsilon]\Gamma_H(W_\varepsilon)
 =\lambda_Hu^2v^2,
 \label{eq:axis-needle}
\end{equation}
where
\[
\begin{array}{c|c|c|c}
H&\text{graph6}&\lambda_H&\kappa_H\\ \hline
137&\texttt{EjsG}&4&1/2\\
139&\texttt{Eju?}&4&1/2\\
145&\texttt{Exe\_}&8&1\\
148&\texttt{EyUG}&4&3/4
\end{array}
\]
Thus the only first-order concentration profiles are the two axes $u=0$ and
$v=0$.  These axes include the exact Turan relocations $(0,1)$ and $(1,0)$,
but also arbitrary partial deletion profiles.

The full arbitrary-fibre interaction on this zero set has an equally simple
form.  For a profile on either axis, let $x\in[0,1]$ denote its nonzero
coordinate.  Fibres belonging to two profiles may be arbitrary measurable
functions with those means.  On the same axis their feasible overlap may be
any value in its full Frechet interval; on opposite axes the relevant overlap
is zero.  The mutual exceptional-cell value $r$ may also be any element of
$[0,1]$.

\begin{theorem}[Arbitrary-fibre axis cone]
\label{thm:axis-cone}
For $H$ equal to Atlas $137,139,145$, or $148$, take any two zero profiles,
on either the same or opposite axes, with nonzero coordinates $x,y$.  For
every feasible fibre overlap and every $r\in[0,1]$, the exact repaired pair
kernel is
\begin{equation}
 K_{H,r}=\kappa_H(1-x)(1-y)\geq0.
 \label{eq:axis-kernel}
\end{equation}
Consequently, for an arbitrary probability distribution of axis profiles and
an arbitrary measurable symmetric exceptional-internal kernel,
\begin{equation}
 [\varepsilon^2]\Gamma_H(W_\varepsilon)
 =\kappa_H\left(\int_S(1-x(z))\,dz\right)^2\geq0.
 \label{eq:axis-integrated}
\end{equation}
The coefficient vanishes only when $x=1$ almost everywhere.  At the level of
connections to the bulk, these are the exact Turan refinement or relocation
profiles; an arbitrary exceptional-internal kernel can still make the full
perturbation nontrivial at higher order.
\end{theorem}

\begin{proof}
Exact enumeration of maps using each of two exceptional roots at most once
gives \eqref{eq:axis-kernel} for $r=0$ and $r=1$, separately for the two
same-axis placements and the opposite-axis placement.  In all $24$ cases the
answer is independent of the shared-neighbour overlap.  The pair kernel is
affine in $r$, so the identity holds on the whole interval.  Substitute
$r=R(z,z')$ in \eqref{eq:integrated-kernel} and apply Fubini:
\[
 \int_{S^2}\kappa_H(1-x(z))(1-x(z'))\,dz\,dz'
 =\kappa_H\left(\int_S(1-x)\right)^2.
\]
Equation \eqref{eq:axis-needle} identifies the asserted zero set.  Finally,
the integral in \eqref{eq:axis-integrated} is zero exactly when $x=1$ almost
everywhere; the qualification about the exceptional-internal kernel follows
because it first appears beyond this quadratic coefficient on that face.
\end{proof}

This theorem closes the entire second-order concentration cone left by the
needle test for these four rows: genuinely two-sided profiles have the strict
first-order margin \eqref{eq:axis-needle}, while every arbitrary mixture on
its zero axes has the nonnegative pair kernel \eqref{eq:axis-kernel}.  It is
still a hierarchical small-support statement.  The independent
two-defect-colouring theorem in
\path{notes/turan_arbitrary_linfty_local_stability.tex} now gives a uniform
$L^\infty$ neighbourhood by a different argument; that theorem controls
uniformly small amplitudes, while the present result permits order-one
changes on a set of vanishing measure.

The only quadratic zero point is the endpoint $x=1$.  Its next coefficient
can also be computed without contracting an arbitrary kernel to its mean.
Split the exceptional set of mass $\varepsilon$ into endpoint types
$S_0,S_1$ of relative masses $\alpha$ and
$\beta=1-\alpha$.  Their bulk profiles are respectively $(0,1)$ and
$(1,0)$.  Let $R_0,R_1$ be arbitrary symmetric graphons inside the two
types, and let $C$ be an arbitrary cross-type kernel.  Write
\[
 r_i=\int R_i,\qquad
 P_i=t(P_3,R_i)=\int d_{R_i}(x)^2\,dx,
 \qquad T_i=t(K_3,R_i).
\]
If $c=\int C$, the raw density is
\begin{equation}
 p_{\rm raw}=\frac12+\varepsilon^2D,
 \qquad
 D=\alpha^2r_0+\beta^2r_1+2\alpha\beta c-2\beta.
 \label{eq:axis-endpoint-density}
\end{equation}
When $D\leq0$, putting the constant value
$-4\varepsilon^2D$ on the bulk $B^2$ cell repairs the density exactly to
$1/2$ and is feasible for all sufficiently small $\varepsilon$.

\begin{theorem}[Arbitrary-kernel axis-endpoint cubic]
\label{thm:axis-endpoint-cubic}
For the endpoint construction above, the raw cubic coefficient and the
cubic coefficient after the formal $B^2$ repair are equal.  They are
independent of the cross kernel $C$ and equal to
\begin{equation}
 [\varepsilon^3]\,t(H,W_\varepsilon)=
 \begin{cases}
 \dfrac{\alpha^3(T_0+2P_0)+\beta^3(T_1+2P_1)}8,
      &H=137,139,\\[6pt]
 \dfrac{\alpha^3P_0+\beta^3P_1}{2},&H=145,\\[6pt]
 \dfrac{\alpha^3P_0+\beta^3P_1}{4},&H=148.
 \end{cases}
 \label{eq:axis-endpoint-cubic}
\end{equation}
In particular this coefficient is nonnegative for arbitrary measurable
$R_0,R_1,C$.  Every fixed endpoint mixture satisfies the desired inequality
for all sufficiently small $\varepsilon$; if $D\leq0$, this holds after the
exact nonnegative repair just described, and if $D>0$, it holds for the raw
graphon.
\end{theorem}

\begin{proof}
At cubic order, direct macro-colouring enumeration has three possible
exceptional vertices.  When they all lie in $S_i$, the surviving internal
selected-edge graphs are one triangle and two copies of the two-edge path
for Atlas $137,139$, two copies of the path for Atlas $145$, and one copy of
the path for Atlas $148$, with the normalization displayed in
\eqref{eq:axis-endpoint-cubic}.  Every mixed-type nonlinear moment has zero
coefficient.  Terms involving only one internal edge depend on
$r_0,r_1,c$; exact expansion of
\eqref{eq:axis-endpoint-density} shows that they cancel.  An edge in the
$B^2$ repair has order $\varepsilon^2$, and its only possible cubic terms
also cancel, proving that the raw and repaired cubics agree.  This finite
coloured-moment enumeration is reconstructed exactly in
\path{codes/needle_higher_order_tests.py}.

All terms on the right of \eqref{eq:axis-endpoint-cubic} are nonnegative.
Suppose first that $D\leq0$.  The repaired graphon has density exactly
$1/2$, so its target is zero and its homomorphism density is nonnegative for
every sufficiently small $\varepsilon$ for which the repair lies in
$[0,1]$.

Suppose now that $D>0$.  The cross term alone cannot cause this inequality:
$2\alpha\beta c\leq2\beta$.  Hence at least one $r_i$ with positive type
mass is nonzero.  Jensen gives $P_i\geq r_i^2>0$, so
\eqref{eq:axis-endpoint-cubic} is strictly positive.  On the other hand,
$p_{\rm raw}-1/2=D\varepsilon^2$; the targets of Atlas $137,139$ vanish to
third order in this density increment and those of Atlas $145,148$ vanish
to second order.  They are therefore respectively $O(\varepsilon^6)$ and
$O(\varepsilon^4)$, while the homomorphism density has a positive
$\varepsilon^3$ coefficient.  The desired inequality follows for all
sufficiently small positive $\varepsilon$.
\end{proof}

The cubic margin can be made uniform, including kernels which depend on
$\varepsilon$.  Indeed, if $D>0$, then
\begin{equation}
 \alpha^2r_0-D-\beta^2
 =\beta^2(1-r_1)+2\alpha\beta(1-c)\geq0.
 \label{eq:axis-endpoint-density-coercivity}
\end{equation}
Thus $\alpha^2r_0\geq D$, and Jensen gives
$P_0\geq r_0^2$.  The cubic functional in
\eqref{eq:axis-endpoint-cubic} is consequently at least $q_HD^2$, where
\[
 q_{137}=q_{139}=q_{148}=\frac14,
 \qquad q_{145}=\frac12.
\]

\begin{corollary}[Uniform arbitrary-kernel endpoint support]
\label{cor:axis-endpoint-uniform}
The endpoint construction is nonnegative uniformly over every choice of
$\alpha,R_0,R_1,C$, even when these data depend on $\varepsilon$, at the
following support radii:
\begin{center}
\begin{tabular}{r|c|c|c}
Atlas&$q_H$&support radius&exact scalar slack\\ \hline
137&$1/4$&$1/10$&$899/15625$\\
139&$1/4$&$1/10$&$899/15625$\\
145&$1/2$&$1/20$&$131441/2000000$\\
148&$1/4$&$1/32$&$3002017/67108864$
\end{tabular}
\end{center}
More precisely, when $D\leq0$ use the exact nonnegative $B^2$ repair to
density $1/2$; when $D>0$, leave the raw graphon unchanged.  In both cases
the desired inequality holds for every positive $\varepsilon$ at most the
displayed radius.
\end{corollary}

\begin{proof}
The terms producing \eqref{eq:axis-endpoint-cubic} come from colourings with
exactly three exceptional vertices.  Retain precisely those nonnegative
terms in the full homomorphism integral.  Every bulk-$A$ mass is
$1/2-\varepsilon=(1/2)(1-2\varepsilon)$, while every bulk-$B$ mass is
$1/2$.  Since $v(H)=6$, their total contribution is at least
\begin{equation}
 t(H,W_{\rm raw})
 \geq \varepsilon^3(1-2\varepsilon)^6q_HD^2.
 \label{eq:axis-endpoint-positive-subintegral}
\end{equation}
This is a positive subintegral, so no signed Taylor remainder is being
discarded.

For Atlas $137,139$, use $0<D\leq1$ and $0\leq p\leq1$ to obtain
\[
 \Phi_H\!\left(\frac12+D\varepsilon^2\right)
 \leq8D^2\varepsilon^6.
\]
For Atlas $145,148$, the factor
$p(3p^2-3p+1)$ is at most one on $[0,1]$, and hence
\[
 \Phi_H\!\left(\frac12+D\varepsilon^2\right)
 \leq4D^2\varepsilon^4.
\]
After removing $D^2\varepsilon^3$, the required scalar margins are
\[
 q_H(1-2\varepsilon)^6-8\varepsilon^3
 \quad(H=137,139),
 \qquad
 q_H(1-2\varepsilon)^6-4\varepsilon
 \quad(H=145,148).
\]
They are decreasing in $\varepsilon$, and at the table radii they equal the
displayed positive rational slacks.  This proves the raw $D>0$ branch.

If $D\leq0$, then $-D\leq2$ and the repair value
$-4\varepsilon^2D$ lies in $[0,1]$ throughout all four displayed ranges.
The repaired density is exactly $1/2$, its target is zero, and its
homomorphism density is nonnegative.  This proves the other branch.
\end{proof}

The same idea closes the profiles which approach the endpoint with the
support scale.  On type $S_0$ let
$F_0:S_0\times B\to[0,1]$ be the arbitrary nonzero fibre, and on type $S_1$
let $F_1:S_1\times A\to[0,1]$ be its counterpart.  In normalized
coordinates put
\[
 d_0=\alpha\iint(1-F_0),\qquad
 d_1=\beta\iint(1-F_1),\qquad d=d_0+d_1.
\]
Retain $r_0,r_1,c$ for the three internal-kernel means and put
\[
 J=\alpha^2r_0+\beta^2r_1+2\alpha\beta c,
 \qquad D_0=J-2\beta.
\]
Direct edge-density integration gives
\begin{equation}
 p_{\rm raw}=\frac12+\varepsilon^2E,
 \qquad
 E=D_0+2d_1-\frac d\varepsilon.
 \label{eq:full-axis-density}
\end{equation}
If $E>0$ and $\varepsilon\leq1/2$, then
\begin{equation}
 E\leq D_0,
 \qquad
 d\leq\frac{\varepsilon D_0}{1-2\varepsilon},
 \qquad
 D_0\leq\alpha^2r_0.
 \label{eq:full-axis-coercivity}
\end{equation}
The first two statements follow by rearranging
\[
 D_0-E=\frac{d_0}{\varepsilon}
       +d_1\left(\frac1\varepsilon-2\right),
\]
and the last is \eqref{eq:axis-endpoint-density-coercivity}.

\begin{theorem}[Uniform arbitrary-fibre full-axis support]
\label{thm:full-axis-uniform}
For Atlas $137,139,145$, or $148$, take an arbitrary measurable distribution
of profiles on the two zero axes, arbitrary exceptional-to-bulk fibres with
those means, and an arbitrary symmetric exceptional-internal kernel.  Split
a set of mass $\varepsilon$ from the first part of $T_2$ as above.  If the
raw density is at least $1/2$, leave the graphon unchanged; otherwise put the
unique constant nonnegative repair on the unchanged bulk $B^2$ cell which
restores density $1/2$.  For every
\[
 0<\varepsilon\leq\frac1{100},
\]
the resulting graphon satisfies the desired inequality.  This radius is
uniform over all profiles and kernels, including data depending on
$\varepsilon$.
\end{theorem}

\begin{proof}
If $E\leq0$, the repair value is $-4\varepsilon^2E$.  Since
$d\leq1$, $J\geq0$, and $2\int_{S_1}x\leq2$, it is at most
$4\varepsilon+8\varepsilon^2<1$ in the asserted range.  The repaired
density is $1/2$, so the target is zero and nonnegativity of the homomorphism
density proves the claim.

Suppose $E>0$.  On the normalized type-$S_0$ space define its rooted fibre
deficit
\[
 \ell(z)=\int_B(1-F_0(z,b))\,db
\]
and take the good set $G_0=\{z:\ell(z)\leq\tau\}$ with
$\tau=1/16$.  Markov's inequality and
\eqref{eq:full-axis-coercivity} give
\[
 \mu(S_0\setminus G_0)
 \leq\frac{d_0}{\alpha\tau}
 \leq\frac{\varepsilon\alpha r_0}
              {(1-2\varepsilon)\tau}.
\]
Deleting the bad vertices removes at most twice their measure from the mean
of $R_0$.  Hence, for
$R_0'=R_0\mathbf1_{G_0\times G_0}$,
\[
 \int R_0'
 \geq r_0\left(1-\frac{32\varepsilon}{1-2\varepsilon}\right).
\]
The two-edge path inequality gives
\[
 t(P_3,R_0')\geq\left(\int R_0'\right)^2.
\]

Retain only the endpoint colourings which produced the $P_0$ part of
\eqref{eq:axis-endpoint-cubic}, and restrict all three exceptional vertices
to $G_0$.  At most eight exceptional-to-bulk fibre factors occur.  For fixed
good exceptional vertices, the union bound
$1-\prod_jf_j\leq\sum_j(1-f_j)$ shows that their integrated product is at
least $1-8\tau=1/2$, even when several factors share a bulk variable.  The
shrinking bulk-$A$ masses cost at most $(1-2\varepsilon)^6$.  Finally,
$\alpha^3r_0^2\geq D_0^2\geq E^2$.  Consequently
\begin{equation}
 t(H,W_{\rm raw})
 \geq q_HE^2\varepsilon^3\,
 \frac12(1-2\varepsilon)^4(1-34\varepsilon)^2,
 \label{eq:full-axis-positive-subintegral}
\end{equation}
where $q_H=1/4,1/4,1/2,1/4$ in Atlas order.

The target bounds used in Corollary~\ref{cor:axis-endpoint-uniform} apply
with $D$ replaced by $E$.  After removing $E^2\varepsilon^3$, the four
scalar margins at $\varepsilon=1/100$ are respectively
\[
 \frac{6276868289}{125000000000},\quad
 \frac{6276868289}{125000000000},\quad
 \frac{3777868289}{62500000000},\quad
 \frac{1277868289}{125000000000}.
\]
Each margin decreases with $\varepsilon$ on $[0,1/100]$, so all are
positive throughout the claimed range.  This proves the raw branch and the
theorem.
\end{proof}

We can now remove the zero-axis restriction completely.  For an arbitrary
profile put
\[
 w(z)=\min\{u(z),v(z)\},\qquad x(z)=\max\{u(z),v(z)\},
\]
and abbreviate
\[
 m=\int_Sw,\qquad n=\int_S(1-x),\qquad
 B_1=\int_Suv,\qquad A_2=\int_Su^2v^2.
\]
Partition $S_0=\{u\leq v\}$ and $S_1=\{v<u\}$, with masses
$\alpha,\beta$, and write
\[
 m_0=\int_{S_0}u,\quad m_1=\int_{S_1}v,\qquad
 d_0=\int_{S_0}(1-v),\quad d_1=\int_{S_1}(1-u).
\]
Thus $m=m_0+m_1$ and $n=d_0+d_1$.  Zeroing the smaller fibre at every
point projects the construction to the full axis cone treated above.  If
$r=\int_{S^2}R$ and $D_0=r-2\beta$, then the same internal-kernel
decomposition as before gives $D_0\leq\alpha^2r_0$.  Direct density
integration gives
\begin{equation}
 \delta:=p_{\rm raw}-\frac12
 =\varepsilon(m-n)
  +\varepsilon^2(D_0+2d_1-2m_0).
 \label{eq:full-profile-density}
\end{equation}

\begin{theorem}[Uniform arbitrary-profile small-support theorem]
\label{thm:full-profile-uniform}
Let $H$ be Atlas $137,139,145$, or $148$.  In the arbitrary-fibre
construction of Section~2, allow every profile $(u(z),v(z))\in[0,1]^2$ and
an arbitrary symmetric exceptional-internal kernel.  If
$p_{\rm raw}\geq1/2$, leave the graphon unchanged.  Otherwise put the unique
constant nonnegative repair on $B^2$ which restores density $1/2$.  For
\[
 0<\varepsilon\leq\frac1{10000},
\]
the resulting graphon satisfies $t(H,W)\geq\Phi_H(p(W))$.  The radius is
uniform over all profiles, fibres, and kernels, including data depending on
$\varepsilon$.
\end{theorem}

\begin{proof}
If $\delta\leq0$, then
$-\delta\leq\varepsilon+2\varepsilon^2$.  The required value on $B^2$ is
$-4\delta\leq4\varepsilon+8\varepsilon^2<1$, so the repair is feasible.
Its density is $1/2$, its target is zero, and $t(H,W)\geq0$ proves the claim.

Suppose $\delta>0$.  Since $d_1\leq n$ and $m_0\geq0$,
\eqref{eq:full-profile-density} implies
\begin{equation}
 (1-2\varepsilon)n<m+\varepsilon D_0.
 \label{eq:full-profile-deficit}
\end{equation}
The exact shifted targets are
\[
\begin{array}{c|c}
H&\Phi_H(1/2+\delta)\\ \hline
137,139&2\delta^3(1+2\delta)^2\\
145,148&\delta^2(1+2\delta)(1+12\delta^2)/2.
\end{array}
\]
For $0\leq\delta\leq1/2$, both are at most $4\delta^2$: the respective
differences from $4\delta^2$ factor as
\begin{equation}
 -2\delta^2(2\delta-1)(2\delta^2+3\delta+2),
 \qquad
 -\frac{\delta^2}{2}(2\delta-1)(12\delta^2+12\delta+7).
 \label{eq:full-profile-target-bound}
\end{equation}

The macro-colourings with exactly one exceptional vertex form a
nonnegative subintegral.  The exact needle identity
\eqref{eq:axis-needle} and independence of the five bulk variables give
\begin{equation}
 t(H,W_{\rm raw})
 \geq c_H\varepsilon(1-2\varepsilon)^5A_2,
 \qquad
 (c_{137},c_{139},c_{145},c_{148})
 =\left(\frac1{16},\frac1{16},\frac18,\frac1{16}\right).
 \label{eq:full-profile-one-vertex}
\end{equation}
No constancy of a fibre is used: conditional on $z$, its incident bulk
variables are integrated independently, giving powers of the fibre means.

First suppose
\begin{equation}
 m\geq2\varepsilon D_0.
 \label{eq:full-profile-one-case}
\end{equation}
Put
\[
 P=\int_S(w+x-1)_+,\qquad N=\int_S(1-w-x)_+.
\]
Pointwise,
\[
 (w+x-1)_+\leq wx=uv,\qquad
 w\leq2wx+(1-w-x)_+.
\]
For the second inequality, if $w+x\geq1$, then $x\geq1/2$ and hence
$w\leq2wx$; when $w+x<1$, its difference is
$(1-x)(1-2w)\geq0$.  Hence
\[
 P\leq B_1,\qquad m\leq2B_1+N.
\]
Also $P-N=m-n$, while positivity of
\eqref{eq:full-profile-density} gives
$N-P<\varepsilon(D_0+2n)$.  Combining these inequalities with
\eqref{eq:full-profile-deficit} yields
\[
 (1-4\varepsilon)m
 <3(1-2\varepsilon)B_1+\varepsilon D_0.
\]
Under \eqref{eq:full-profile-one-case} and
$\varepsilon\leq1/100$, this implies $m\leq7B_1$.  By
Cauchy--Schwarz, $B_1^2\leq A_2$.  A second use of
\eqref{eq:full-profile-deficit} gives
\[
 \delta
 \leq\varepsilon m+\varepsilon^2(D_0+2n)
 \leq\frac{3\varepsilon m}{2(1-2\varepsilon)}
 \leq\frac{75}{49}\varepsilon m.
\]
Consequently
\[
 \Phi_H(1/2+\delta)
 \leq4\delta^2
 \leq\frac{22500}{49}\varepsilon^2A_2.
\]
After removing $\varepsilon A_2$ from
\eqref{eq:full-profile-one-vertex}, the smallest scalar margin at
$\varepsilon=1/10000$ is
\begin{equation}
 \frac1{16}(1-2\varepsilon)^5-\frac{22500}{49}\varepsilon
 =\frac{40471936237751224951}
        {2450000000000000000000}>0.
 \label{eq:full-profile-one-slack}
\end{equation}
The first term decreases and the subtracted term increases with
$\varepsilon$, so the margin is positive throughout the asserted interval.
If $A_2=0$, the same density relations make this case incompatible with
$\delta>0$; equivalently the bound follows by continuity without division
by $A_2$.

It remains to suppose $m<2\varepsilon D_0$, so $D_0>0$.  Equation
\eqref{eq:full-profile-deficit} gives
\begin{equation}
 n<\frac{3\varepsilon D_0}{1-2\varepsilon}.
 \label{eq:full-profile-internal-deficit}
\end{equation}
Project each profile to its nearer axis by deleting its smaller fibre.  The
original graphon dominates this projection pointwise.  On $S_0$, take the
good set where the rooted deficit of the remaining $B$-fibre is at most
$1/16$.  By \eqref{eq:full-profile-internal-deficit}, Markov's inequality,
and $D_0\leq\alpha^2r_0$, its relative bad-set measure is at most
\[
 \frac{48\varepsilon\alpha r_0}{1-2\varepsilon}.
\]
Deleting bad endpoints from the internal kernel therefore leaves edge mean
at least
\[
 r_0\left(1-\frac{96\varepsilon}{1-2\varepsilon}\right)
 =r_0\frac{1-98\varepsilon}{1-2\varepsilon}.
\]
The same two-edge-path inequality and eight-factor union bound used in
Theorem~\ref{thm:full-axis-uniform} now give
\begin{equation}
 t(H,W_{\rm raw})\geq
 q_HD_0^2\varepsilon^3
 \frac12(1-2\varepsilon)^4(1-98\varepsilon)^2,
 \qquad
 (q_{137},q_{139},q_{145},q_{148})
 =\left(\frac14,\frac14,\frac12,\frac14\right).
 \label{eq:full-profile-internal}
\end{equation}
Here $\alpha^3r_0^2\geq D_0^2$, because
$D_0\leq\alpha^2r_0$ and $\alpha\leq1$.
Only colourings compatible with the projected axis graphon were retained,
so pointwise domination justifies this subintegral in the original graphon.

Finally, \eqref{eq:full-profile-density},
\eqref{eq:full-profile-internal-deficit}, and
$m<2\varepsilon D_0$ give
\[
 \delta<\frac{3\varepsilon^2D_0}{1-2\varepsilon}.
\]
Thus the target is at most
$36\varepsilon^4D_0^2/(1-2\varepsilon)^2$.  After removing
$\varepsilon^3D_0^2$ from \eqref{eq:full-profile-internal}, the smallest
scalar margin at $\varepsilon=1/10000$ is
\begin{equation}
 \frac18(1-2\varepsilon)^4(1-98\varepsilon)^2
 -\frac{36\varepsilon}{(1-2\varepsilon)^2}
 =\frac{371296887848737563915027482401}
        {3123750125000000000000000000000}>0.
 \label{eq:full-profile-internal-slack}
\end{equation}
Again the first term decreases and the subtracted term increases, so the
margin is positive on the whole interval.  This proves the second case.
\end{proof}

Thus the complete single-exceptional-set small-support model is settled for
these four rows, including arbitrary two-sided profiles approaching either
zero axis at any rate.  This remains a fixed-representative hierarchical
statement and is not a cut-distance theorem.

\begin{remark}[Higher-order mean contraction is invalid]
At the endpoint $x=1$, the exceptional-internal kernel is invisible to
\eqref{eq:axis-integrated}, but it cannot be replaced by its mean at cubic
order.  Give each endpoint type mass $\varepsilon/2$, split the first type
into two equal atoms, set the second-type and cross-type kernels to zero, and
put the exact repair $7\varepsilon^2/2$ on $B^2$.  The two first-type kernels
\[
R_{\rm const}\equiv\frac12,\qquad
 R_{\rm var}=\begin{pmatrix}1&1/2\\1/2&0\end{pmatrix}
\]
have the same mean.  Nevertheless their exact cubic coefficients are
\[
\begin{array}{c|cc}
H&R_{\rm const}&R_{\rm var}\\ \hline
137,139&5/512&27/2048\\
145&1/64&5/256\\
148&1/128&5/512.
\end{array}
\]
Thus a higher-order proof must retain rooted-degree and mixed internal
moments.  This is a failed contraction, not a negative direction; all eight
displayed coefficients are positive.  The full warning is recorded as
Attempt 28 in \path{FAILED_PROOF_ATTEMPTS.md}.
\end{remark}

\section{Exact boundary: Atlas 152}

The same pointwise method is false for the sixth flat row, Atlas 152.
For mutual cell $r=0$, profiles
\[
 (u,v)=\left(\frac16,1\right),\qquad
 (u',v')=\left(\frac56,\frac23\right)
\]
and indicator fibres with $h_0=0$, $h_1=2/3$ give exactly
\[
 K_{152,0}(z,z')=-\frac1{324}.
\]
This is only an off-diagonal failure.  The full two-profile matrix is
\[
 \begin{pmatrix}5/288&-1/324\\-1/324&2/9\end{pmatrix},
 \qquad \det=\frac{101}{26244}>0.
\]
Thus it is not a negative quadratic direction.  Atlas 152 requires an
integrated copositivity or Gram-moment inequality; pointwise pair-kernel
nonnegativity cannot prove it.

For reference, the hard integrated inequality has a compact exact
Hilbert-space form.  Write $\mathbb E$ for integration over the normalized
exceptional space and set
\[
\begin{array}{ll}
 P(a,a')=\mathbb E[vF(a)F(a')],&
 Q(b,b')=\mathbb E[uG(b)G(b')],\\
 X(a)=\mathbb E[uvF(a)],&Y(a)=\mathbb E[vF(a)],\\
 Z(b)=\mathbb E[uvG(b)],&T(b)=\mathbb E[uG(b)],\\
 L(a,b)=\mathbb E[uF(a)G(b)],&
 M(a,b)=\mathbb E[vF(a)G(b)],
\end{array}
\]
and
\[
 C=\mathbb E[uv],\qquad
 \delta=\mathbb E[u+v-1].
\]
Direct integration of the exact pair kernel at the worst mutual value
$R=0$ gives
\begin{equation}
 8[\varepsilon^2]\Gamma_{152}(W_\varepsilon)
 =\|P\|_2^2+\|Q\|_2^2
 +2\langle X,Y\rangle+2\langle Z,T\rangle
 +2\langle L,M\rangle+C^2-4\delta^2.
 \label{eq:atlas152-hilbert}
\end{equation}
An arbitrary internal kernel only adds
$\frac14\int R(z,z')u(z)v(z)u(z')v(z')\,dz\,dz'\geq0$ to the quadratic
coefficient, so \eqref{eq:atlas152-hilbert} is the exact hard case.

There is a useful joint factorization which shows exactly what geometry is
still missing.  Put
\[
 h_F(z,z')=\int_A F_zF_{z'},\qquad
 h_G(z,z')=\int_B G_zG_{z'}
\]
and define the generally nonsymmetric nonnegative kernel
\[
 \mathcal A(z,z')=
 v(z)h_F(z,z')+u(z)h_G(z,z')+u(z)v(z).
\]
Expanding the product gives the exact identity
\begin{equation}
 \|P\|_2^2+\|Q\|_2^2+2\langle X,Y\rangle
 +2\langle Z,T\rangle+2\langle L,M\rangle+C^2
 =\int_{S^2}\mathcal A(z,z')\mathcal A(z',z)\,dz\,dz'.
 \label{eq:joint-factor}
\end{equation}
Thus the positive terms are one structured joint kernel, not six unrelated
moments.

A tempting completion of this observation is nevertheless false.  With
$A=\mathbb E[u^2v]$ and $B=\mathbb E[uv^2]$, the integral of
$\mathcal A$ is $A+B+C$.  Moreover
\[
 uv(u+v+1)\geq2(u+v-1)\qquad (u+v\geq1).
\]
For example, after homogenizing on the triangle with vertices
$(1,0),(0,1),(1,1)$, the six degree-three barycentric coefficients of the
difference are $(2,2,3,1,1,1)$.  Hence $A+B+C\geq2\delta$.  It would
therefore suffice to prove
\begin{equation}
 \int\mathcal A(z,z')\mathcal A(z',z)\,dz\,dz'
 \geq(A+B+C)^2.
 \label{eq:false-joint-contraction}
\end{equation}
But \eqref{eq:false-joint-contraction} is not a valid inequality even for
the structured kernel above.  Give $S$ two atoms of mass $1/2$, $A$ two
atoms of mass $1/2$, and $B$ one atom.  Take
\[
 F_1=(0,1/3),\quad F_2=(1,1/3),\qquad
 G_1=1,\quad G_2=1/3.
\]
Then
\[
 (u_1,v_1)=(1/6,1),\qquad (u_2,v_2)=(2/3,1/3),
\]
so both profiles lie in $\Delta_+$.  Exact evaluation gives
\[
 A=\frac{19}{216},\quad B=\frac{13}{108},\quad C=\frac7{36},
 \quad \delta=\frac1{12},
\]
and
\[
 \int\mathcal A(z,z')\mathcal A(z',z)=\frac{1867}{11664},
 \qquad
 \frac{1867}{11664}-(A+B+C)^2=-\frac{101}{46656}.
\]
This does not give a negative Atlas 152 direction: the true remainder is
\[
 \frac{1867}{11664}-4\delta^2=\frac{1543}{11664}>0.
\]
It does show that even retaining the joint factorization is insufficient if
one then applies an unrestricted ``trace versus total mass'' contraction.
The required lower bound must depend directly on the density displacement
$\delta$, and must exploit more than the total mass of $\mathcal A$.

A pointwise symmetrization of the two orientations also loses essential
information.  Indeed, it would be enough to prove
\begin{equation}
 \int\sqrt{\mathcal A(z,z')\mathcal A(z',z)}\,dz\,dz'
 \geq2\delta,
 \label{eq:false-geometric-contraction}
\end{equation}
because Cauchy--Schwarz would then give the desired lower bound for
$\int\mathcal A\mathcal A^T$.  But
\eqref{eq:false-geometric-contraction} is false even for exact admissible
fibres.  Give $S$, $A$, and $B$ two equal atoms each and take
\[
 F_1=(0,0),\quad F_2=(0,1/2),\qquad
 G_1=(1,1),\quad G_2=(1,1).
\]
The mean profiles are $(0,1)$ and $(1/4,1)$, and
\[
 \bigl(\mathcal A(i,j)\bigr)_{i,j=1}^2
 =\begin{pmatrix}0&0\\[2pt]1/2&5/8\end{pmatrix},
 \qquad \delta=\frac18.
\]
Thus the left side of \eqref{eq:false-geometric-contraction} is $5/32$,
short of $2\delta=1/4$ by $3/32$.  Nevertheless the actual positive joint
integral is $25/256$, so the true eightfold remainder is
\[
 \frac{25}{256}-4\delta^2=\frac9{256}>0.
\]
This second obstruction rules out a symmetrized scalar contraction, not the
structured Atlas 152 functional itself.

Nor may the overlap variables be minimized independently.  For the two
profiles
\[
 (u_1,v_1)=(49/50,2/5),\qquad (u_2,v_2)=(1/50,1),
\]
substituting the individually sharp diagonal bounds
$h_F(i,i)\geq u_i^2$, $h_G(i,i)\geq v_i^2$ and the off-diagonal Frechet
bounds produces
\[
 M_{\rm rel}=
 \begin{pmatrix}
 11438061/312500000&-33/31250\\
 -33/31250&201/50000000
 \end{pmatrix},
 \qquad
 \det M_{\rm rel}=-\frac{15124949739}{15625000000000000}<0.
\]
This relaxed matrix is not copositive.  However, the bounds used in it are
not jointly realizable.  Complementary indicator $F$ fibres attain the
off-diagonal overlap zero, but then their diagonal overlaps are $u_i$, not
$u_i^2$.  Together with constant $G$ fibres, the resulting compatible
matrix is
\[
 M_{\rm comp}=
 \begin{pmatrix}
 120119/3125000&-33/31250\\
 -33/31250&1/4000
 \end{pmatrix},
 \qquad
 \det M_{\rm comp}=\frac{530899}{62500000000}>0.
\]
Thus this is not a negative direction.  It proves instead that common-fibre
Gram consistency is essential and cannot be replaced by independent
pairwise overlap intervals.

\subsection{Every two-profile indicator mixture is nonnegative}

The preceding compatible example is a special case of a complete exact
two-profile theorem.  Suppose the exceptional space has at most two profile
classes and, in each class, both fibres are indicators of measurable sets.
Write their mean profiles as
\[
 (u,v),(u',v')\in\Delta_+,
\]
their mixed intersections as $h_F,h_G$, and let $K$ be the mutual-cell-zero
pair kernel.  On a diagonal, the indicator identities
$h_F=u$, $h_G=v$ give
\begin{equation}
 D(u,v)=
 \frac{9u^2v^2-4(u+v-1)^2}{8}
 =\frac{(3uv-2u-2v+2)(3uv+2u+2v-2)}8.
 \label{eq:indicator-diagonal}
\end{equation}
Since $uv\geq u+v-1\geq0$, one has
$D(u,v)\geq5(u+v-1)^2/8\geq0$.

Every positive-degree coefficient of the mixed kernel in $h_F,h_G$ is
nonnegative.  Hence its smallest feasible value $K_{\min}$ is obtained at
\[
 h_F=\max\{0,u+u'-1\},\qquad
 h_G=\max\{0,v+v'-1\}.
\]
The two Fr\'echet switches cut $\Delta_+^2$ into the same six exact
four-simplices used above.  On each simplex, pull back the degree-eight
determinant polynomial
\begin{equation}
 D(u,v)D(u',v')-K_{\min}(u,v,u',v')^2.
 \label{eq:indicator-determinant}
\end{equation}
After homogeneous barycentric substitution, every coefficient is
nonnegative without degree elevation.  The exact signatures
\[
\begin{array}{c|ccc}
\text{simplex}&\text{degree}&\text{nonzero coefficients}&
 \text{least positive coefficient}\\ \hline
1&8&365&9/64\\
2&8&365&9/64\\
3&8&410&9/64\\
4&8&474&9/64\\
5&8&244&21/64\\
6&8&244&21/64
\end{array}
\]
are independent exact certificates of \eqref{eq:indicator-determinant}.

\begin{theorem}[Atlas 152, two indicator-fibre profiles]
\label{thm:atlas152-two-indicator}
On the flat threshold face of Atlas $152$, the complete quadratic
coefficient is nonnegative for every mixture supported on at most two
profiles whose two exceptional-to-bulk fibres are indicators.  The two
profiles, their masses, their set intersections, and the
exceptional--exceptional kernel are arbitrary subject only to feasibility.
\end{theorem}

\begin{proof}
Let the two masses be $s,t\geq0$.  If the actual mixed entry $K$ is
nonnegative, then
\[
 s^2D(u,v)+2stK+t^2D(u',v')\geq0.
\]
If $K<0$, overlap monotonicity gives
$K_{\min}\leq K<0$, hence $K^2\leq K_{\min}^2$.
Equation~\eqref{eq:indicator-determinant} then makes the two-by-two matrix
positive semidefinite.  Finally, a nonzero exceptional--exceptional kernel
adds
\[
 \frac14\int R(z,z')u(z)v(z)u(z')v(z')\,dz\,dz'\geq0.
\]
\end{proof}

This theorem supplies directions outside both the coarse-profile and
comonotone cones.  In particular, the two indicator sets in either bulk part
may cross arbitrarily.  Thus any unresolved Atlas $152$ quadratic direction
must use at least three profile classes or a genuinely fractional fibre in
at least one of the two bulk parts.

\subsection{Arbitrary mixtures on the two profile axes}

There is also a profile-count-free theorem.  First observe that profiles of
zero density displacement may always be discarded.  Indeed, put
$d(z)=u(z)+v(z)-1$ and let $E_+=\{d>0\}$ have mass $\alpha$.  If
$\alpha=0$, then $\delta=0$ and the unpaired joint-kernel integral is already
nonnegative.  Otherwise, normalize the measure on $E_+$.  Since
$\mathcal A(z,z')\mathcal A(z',z)\geq0$, deleting every pair meeting
$E_+^c$ can only lower that integral, while
$\delta=\alpha\delta_+$.  Thus any lower bound on the normalized positive
part scales by the same factor $\alpha^2$ as the required $4\delta^2$.

Now suppose each positive-displacement profile lies on
\[
 \{u=1\}\cup\{v=1\}.
\]
If $u=1$, then its first fibre equals one almost everywhere; if $v=1$, the
same is true of its second fibre.  For two profiles $(1,x),(1,y)$ on the
first axis, write $h\geq0$ for the overlap of their second fibres.  Direct
substitution into the exact mutual-cell-zero pair kernel gives
\begin{equation}
 K_{152,0}\bigl((1,x),(1,y)\bigr)
 =\frac{h(h+2x+2y)}8\geq0.
 \label{eq:atlas152-same-axis}
\end{equation}
The other axis is symmetric.  Across the axes, for profiles $(1,x)$ and
$(y,1)$, the forced overlaps are $h_F=y$ and $h_G=x$, and hence
\begin{equation}
 K_{152,0}\bigl((1,x),(y,1)\bigr)
 =\frac{xy(xy+2x+2y)}8\geq0.
 \label{eq:atlas152-cross-axis}
\end{equation}

\begin{theorem}[Atlas 152, arbitrary axis-profile mixtures]
\label{thm:atlas152-axis}
On the flat threshold face of Atlas $152$, the complete quadratic
coefficient is nonnegative whenever almost every positive-displacement
profile satisfies $u=1$ or $v=1$.  The profile distribution, the number of
profile classes, the unconstrained fibre on each axis, all fibre crossings,
and the exceptional--exceptional kernel are arbitrary.
\end{theorem}

\begin{proof}
Discard zero-displacement profiles by the preceding scaling argument.
Equations~\eqref{eq:atlas152-same-axis} and
\eqref{eq:atlas152-cross-axis} make the mutual-cell-zero pair kernel
pointwise nonnegative on the remaining support.  Integrate it.  As before,
the exceptional--exceptional kernel contributes only a nonnegative term.
\end{proof}

Consequently a still-unresolved indicator construction must not only use at
least three crossing profiles; a set of positive exceptional mass must also
carry strict-interior profiles
\[
 u<1,\qquad v<1,\qquad u+v>1.
\]

\subsection{Arbitrary fibres over a chain of mean profiles}

The profile-count restriction can be removed in a different direction.  Say
that the mean-profile support is a \emph{coordinatewise chain} if almost every
pair is comparable in the product order.  By symmetry, write a comparable
pair as
\[
 u\leq u',\qquad v\leq v'.
\]
Every positive-degree coefficient of the Atlas $152$ pair kernel in
$h_F,h_G$ is nonnegative.  It therefore suffices to use the two sharp lower
overlaps.

There is a compact exact certificate on the complete comparable-profile
domain.  Introduce deficit and increment coordinates
\[
 u=1-d-y,\qquad v=1-c-x,
 \qquad u'=1-d,\qquad v'=1-c.
\]
The domain is the standard four-simplex
\begin{equation}
 c,d,x,y\geq0,\qquad c+d+x+y\leq1,
 \label{eq:comparable-simplex}
\end{equation}
and the lower overlaps are
\[
 h_F=(1-2d-y)_+,\qquad h_G=(1-2c-x)_+.
\]
The region where both displayed affine forms are nonnegative has the eight
vertices
\[
\begin{array}{c|c@{\qquad}c|c@{\qquad}c}
0&(0,0,0,0)&1&(0,0,0,1)\\
2&(0,0,1,0)&3&(0,1/2,0,0)\\
4&(0,1/2,1/2,0)&5&(1/2,0,0,0)\\
6&(1/2,0,0,1/2)&7&(1/2,1/2,0,0).
\end{array}
\]
It is triangulated by the seven four-simplices
\[
13467,\quad01246,\quad01346,\quad24567,\quad02456,\quad34567,\quad03456,
\]
where a string lists its five vertex labels.  The side on which $h_F$ is
positive and $h_G=0$ is the simplex with vertices
\[
 (0,0,1,0),(1,0,0,0),(1/2,0,0,0),
 (1/2,0,0,1/2),(1/2,1/2,0,0),
\]
and the other side is obtained by $(c,d,x,y)\mapsto(d,c,y,x)$.  The region
where both lower overlaps vanish is only the boundary point
$(1/2,1/2,0,0)$.

For completeness, the base simplex has volume $1/24$, the two side
simplices have volume $1/192$ each, and the seven central simplex volumes
sum to $1/32$; hence they sum to $1/24$.  There is also an exact
face-to-face certificate.  Convert each simplex to its five rational
barycentric half-spaces, intersect every pair by enumerating four active
facets, and compare every resulting intersection vertex with the convex
hull of the vertices shared by that pair.  In all $36$ pairs the two sets
agree.  Classified by the number of shared vertices, the incidence counts
are
\[
 1:3,\qquad 2:9,\qquad 3:13,\qquad 4:11.
\]
Thus all relative interiors are disjoint.  Since every simplex lies in
\eqref{eq:comparable-simplex} and their exact volumes sum to its volume,
they form a partition; the remaining faces follow by continuity.

On each simplex substitute the appropriate lower-overlap formula into
$K_{152,0}$ and homogenize in its five barycentric coordinates.  Every
coefficient is nonnegative without degree elevation.  In the order of the
seven central simplices followed by the two sides, the exact signatures are
\[
\begin{array}{c|ccc}
&\text{degree}&\text{nonzero coefficients}&\text{least positive coefficient}\\ \hline
13467&4&65&1/128\\
01246&4&68&3/64\\
01346&4&65&1/32\\
24567&4&65&1/128\\
02456&4&65&1/32\\
34567&4&68&1/128\\
03456&4&68&1/32\\
\text{first side}&4&53&1/128\\
\text{second side}&4&53&1/128
\end{array}
\]
At the remaining boundary point the value is $1/128$.

\begin{theorem}[Atlas 152, comparable mean profiles]
\label{thm:atlas152-comparable-means}
On the flat threshold face of Atlas $152$, the complete quadratic
coefficient is nonnegative whenever the mean profiles form a coordinatewise
chain.  The profile distribution and the fibres themselves are arbitrary;
in particular, the fibres may be fractional and may cross without any
nesting assumption.
\end{theorem}

\begin{proof}
For almost every pair of exceptional points, orient the two comparable mean
profiles as above.  Overlap monotonicity and the nine exact certificates give
$K_{152,0}(z,z')\geq0$.  Integrate this pointwise inequality.  The
exceptional--exceptional correction is nonnegative.
\end{proof}

Thus a residual Atlas $152$ obstruction must give positive product measure
to incomparable profile pairs, equivalently pairs for which
$(u-u')(v-v')<0$.

\subsection{Arbitrary fibres at one density displacement}

A second pointwise theorem covers the transverse antichain direction.  Fix
$t\in[0,1]$ and suppose every mean profile satisfies
\[
 u+v-1=t.
\]
For a pair, assume $u\leq u'$ and write
\[
 u=t+a,\qquad v=1-a,\qquad
 u'=t+a+b,\qquad v'=1-a-b.
\]
The full domain is the standard three-simplex
\begin{equation}
 t,a,b\geq0,\qquad t+a+b\leq1,
 \label{eq:equal-displacement-simplex}
\end{equation}
and the lower overlaps are
\[
 h_F=(2t+2a+b-1)_+,\qquad
 h_G=(1-2a-b)_+.
\]
The two switch planes partition \eqref{eq:equal-displacement-simplex} into
four tetrahedra.  In order from the $h_F=0$ side through the two central
tetrahedra to the $h_G=0$ side, their vertices are
\[
\begin{array}{c|l}
1&(0,0,0),(0,0,1),(0,1/2,0),(1/2,0,0)\\
2&(0,0,1),(1,0,0),(1/2,0,0),(1/2,1/2,0)\\
3&(0,0,1),(0,1/2,0),(1/2,0,0),(1/2,1/2,0)\\
4&(0,0,1),(0,1,0),(0,1/2,0),(1/2,1/2,0).
\end{array}
\]
Each has exact volume $1/24$.  Exact rational facet enumeration checks all
six pairwise intersections: three pairs share exactly two vertices and
three share exactly three, and in every case the intersection is precisely
the convex hull of those common vertices.  Hence the interiors are disjoint,
and their total volume $1/6$ proves that they fill the base simplex.
Pulling back the appropriate lower-envelope pair kernel gives the exact
coefficient signatures
\[
 (4,22,1/128),\quad(4,26,1/16),\quad
 (4,27,1/128),\quad(4,22,1/128).
\]
Thus every coefficient is nonnegative without degree elevation.

\begin{theorem}[Atlas 152, equal displacement]
\label{thm:atlas152-equal-displacement}
On the flat threshold face of Atlas $152$, the complete quadratic
coefficient is nonnegative whenever $u(z)+v(z)-1$ is almost everywhere
constant.  The profile distribution and the fibres are arbitrary; in
particular, the mean profiles may traverse the complete decreasing line
segment and the fractional fibres may cross arbitrarily.
\end{theorem}

\begin{proof}
Orient each pair by its first coordinate.  Overlap monotonicity and the four
simplex certificates prove its mutual-cell-zero kernel nonnegative.  Integrate
and add the nonnegative exceptional--exceptional correction.
\end{proof}

Consequently a residual obstruction must combine incomparable mean profiles
with genuinely varying density displacement.

\subsection{Arbitrary fibres over one mean profile}

Fractional fibres are also harmless when their two means do not vary with
the exceptional point.  Suppose $u(z)=u$ and $v(z)=v$ almost everywhere,
while the functions $F_z$ and $G_z$ are otherwise arbitrary.  Then the joint
kernel is symmetric:
\[
 \mathcal A(z,z')=v h_F(z,z')+u h_G(z,z')+uv
                 =\mathcal A(z',z).
\]
Moreover, with $\overline F=\mathbb E_zF_z$ and
$\overline G=\mathbb E_zG_z$,
\[
 \mathbb E_{z,z'}h_F(z,z')=\|\overline F\|_2^2\geq u^2,
 \qquad
 \mathbb E_{z,z'}h_G(z,z')=\|\overline G\|_2^2\geq v^2.
\]
Writing $d=u+v-1$, Jensen therefore gives
\begin{align*}
 \int\mathcal A(z,z')\mathcal A(z',z)\,dz\,dz'
 &=\mathbb E\mathcal A^2\\
 &\geq(\mathbb E\mathcal A)^2\\
 &\geq\bigl(uv(u+v+1)\bigr)^2\\
 &\geq4d^2,
\end{align*}
where the final unsquared slack is exactly
\[
 uv(u+v+1)-2d
 =uvd+2(1-u)(1-v)\geq0.
\]

\begin{theorem}[Atlas 152, one mean profile with arbitrary fibres]
\label{thm:atlas152-single-mean}
On the flat threshold face of Atlas $152$, the complete quadratic
coefficient is nonnegative when the exceptional-to-bulk fibre means are one
fixed profile $(u,v)\in\Delta_+$.  The individual fibres may be arbitrary
$[0,1]$-valued measurable functions and may vary and cross without
restriction over the exceptional space.
\end{theorem}

The exceptional--exceptional correction is again nonnegative, so the
preceding calculation proves the theorem.  In particular, a residual
fractional-fibre construction must have nonzero mean-profile dispersion.

The argument is quantitatively stable under dispersion of the mean profile.
Let
\[
 \bar u=\mathbb E u(z),\qquad
 \bar v=\mathbb E v(z),\qquad
 \bar d=\bar u+\bar v-1,
\]
and define the exact fixed-profile margin
\begin{equation}
 \mathcal G(\bar u,\bar v)
 =\bigl(\bar u\bar v(\bar u+\bar v+1)\bigr)^2-4\bar d^2.
 \label{eq:atlas152-profile-margin}
\end{equation}
It is nonnegative on $\Delta_+$, because it is the product of the
nonnegative sum and difference whose difference factor is
\[
 \bar u\bar v\bar d+2(1-\bar u)(1-\bar v).
\]
It vanishes only at $(\bar u,\bar v)=(1,0),(0,1)$.

Keep the actual overlap kernels but freeze only the three scalar
coefficients:
\[
 \mathcal A_0(z,z')
 =\bar v h_F(z,z')+\bar u h_G(z,z')+\bar u\bar v.
\]
This kernel is symmetric.  Put
\[
 e_u(z)=u(z)-\bar u,
 \qquad
 e_v(z)=v(z)-\bar v.
\]
Since every scalar and overlap lies in $[0,1]$,
\[
 |\mathcal A(z,z')-\mathcal A_0(z,z')|
 \leq2\bigl(|e_u(z)|+|e_v(z)|\bigr),
 \qquad 0\leq\mathcal A,\mathcal A_0\leq3.
\]
Applying this estimate in the two orientations and integrating yields
\begin{equation}
 \int\mathcal A(z,z')\mathcal A(z',z)
 \geq\int\mathcal A_0(z,z')^2
 -12\bigl(\mathbb E|e_u|+\mathbb E|e_v|\bigr).
 \label{eq:atlas152-dispersion-error}
\end{equation}
The same Gram-marginal and Jensen calculation as above gives
\[
 \int\mathcal A_0^2
 \geq\bigl(\bar u\bar v(\bar u+\bar v+1)\bigr)^2.
\]
Since the aggregate displacement is exactly $\bar d$, we obtain the
following uniform statement.

\begin{corollary}[Atlas 152, dispersed mean profiles]
\label{cor:atlas152-profile-concentration}
Let
\[
 \sigma_u=\bigl(\mathbb E e_u^2\bigr)^{1/2},
 \qquad
 \sigma_v=\bigl(\mathbb E e_v^2\bigr)^{1/2}.
\]
The complete quadratic coefficient is nonnegative whenever
\[
 12\bigl(\mathbb E|e_u|+\mathbb E|e_v|\bigr)
 \leq\mathcal G(\bar u,\bar v).
\]
The fibres themselves remain arbitrary measurable $[0,1]$-valued functions.
By Cauchy--Schwarz, the simpler sufficient condition
\[
 12(\sigma_u+\sigma_v)\leq\mathcal G(\bar u,\bar v)
\]
also applies.  In particular, if
\begin{equation}
 |u(z)-\bar u|\leq\varepsilon,
 \qquad
 |v(z)-\bar v|\leq\varepsilon
 \quad\text{almost everywhere},
 \label{eq:atlas152-profile-concentration}
\end{equation}
then it suffices that
\[
 \varepsilon\leq\frac{\mathcal G(\bar u,\bar v)}{24}.
\]
Thus every average profile other than $(1,0),(0,1)$ has explicit positive
$L^1$-dispersion, $L^2$-dispersion, and $L^\infty$-concentration
neighbourhoods containing no Atlas $152$ quadratic obstruction.
\end{corollary}

\subsection{A stronger $C$-contraction is false}

The joint representation and the elementary inequality
$uv\geq u+v-1$ suggest trying to prove
\[
 \int\mathcal A(z,z')\mathcal A(z',z)\,dz\,dz'
 \stackrel{?}{\geq}4C^2,
 \qquad C=\mathbb E[uv].
\]
This would be stronger than the needed $4\delta^2$ bound.  It is false even
on the boundary $\delta=0$ with two exact compatible profiles.

Give the two exceptional classes masses $3/5,2/5$.  On the first bulk
space take constant fibres
\[
 F_1\equiv\frac9{10},\qquad F_2\equiv\frac14.
\]
On a four-atom second bulk space take
\[
 G_1=(2/5,0,0,0),\qquad G_2=(0,1,1,1).
\]
The mean profiles are $(9/10,1/10)$ and $(1/4,3/4)$, so each has zero
density displacement.  The exact joint-kernel matrix is
\[
 \bigl(\mathcal A(i,j)\bigr)_{i,j=1}^2
 =
 \begin{pmatrix}
 207/1000&9/80\\
 57/160&27/64
 \end{pmatrix}.
\]
Consequently
\[
 C=\frac{129}{1000},\qquad
 \int\mathcal A\mathcal A^T=\frac{25255881}{400000000},
\]
and
\[
 \frac{25255881}{400000000}-4C^2
 =-\frac{1369719}{400000000}<0.
\]
The true Atlas $152$ remainder is nevertheless
$25255881/400000000>0$, because $\delta=0$.  Thus replacing the density
displacement by the larger moment $C$ loses precisely the boundary geometry
that the desired inequality may exploit.

There is a broad compatible subcone on which the desired pointwise estimate
does hold.  Suppose both fibre families are comonotone: after
measure-preserving parametrizations one may write
\[
 F_z(a)=\mathbf1_{a\leq u(z)},
 \qquad
 G_z(b)=\mathbf1_{b\leq v(z)}.
\]
Then
\[
 h_F(z,z')=\min\{u,u'\},\qquad
 h_G(z,z')=\min\{v,v'\}.
\]
For fixed $(u,v)\in\Delta_+$, the function
\[
 (u',v')\longmapsto
 v\min\{u,u'\}+u\min\{v,v'\}
\]
is concave on the upper profile triangle.  Its minimum occurs at a vertex.
At the three vertices, its values are respectively
\[
 uv,\qquad uv,\qquad 2uv,
\]
and
\[
 uv-(u+v-1)=(1-u)(1-v)\geq0.
\]
Consequently
\[
 \mathcal A(z,z')
 \geq uv+(u+v-1)\geq2(u+v-1).
\]
The same argument in the reverse orientation gives
$\mathcal A(z',z)\geq2(u'+v'-1)$.  Therefore the Atlas $152$ pair kernel is
pointwise nonnegative throughout the comonotone-fibre cone.  Together with
Theorems~\ref{thm:atlas152-two-indicator}, \ref{thm:atlas152-axis},
\ref{thm:atlas152-comparable-means},
\ref{thm:atlas152-equal-displacement}, and~\ref{thm:atlas152-single-mean},
this leaves only mixtures with a positive mass of strict-interior profiles,
sufficiently large mean-profile dispersion, a positive measure of
incomparable profile pairs, and genuinely varying density displacement; in
the indicator case at least three profiles must cross.

The most immediate contraction estimate does not suffice.  If
$A=\mathbb E[u^2v]$ and $B=\mathbb E[uv^2]$, then $P$ has marginal $X$,
$X\leq Y$, and their $B$-analogues give
\[
 \|P\|_2^2+2\langle X,Y\rangle\geq3A^2,
 \qquad
 \|Q\|_2^2+2\langle Z,T\rangle\geq3B^2.
\]
Also, with $J=\mathbb E[FG]$, one has $0\leq L,M\leq J$ and
$L+M\geq J$.  Pointwise normalization shows
$LM\geq(L+M-J)^2$, hence
\[
 \langle L,M\rangle\geq(A+B-C)^2.
\]
The resulting scalar lower bound is
\[
 3A^2+3B^2+2(A+B-C)^2+C^2-4\delta^2.
\]
It is not nonnegative: for equal masses on profiles
$(1,\eta)$ and $(\eta,1)$ it equals
\[
 \eta^2\left(-\frac32+3\eta+\frac72\eta^2\right),
\]
which is negative for small positive $\eta$.  Therefore any successful proof
of \eqref{eq:atlas152-hilbert} must retain more of the joint geometry of the
six functions; separate marginal contractions lose essential information.

\section{The residual cone contains an exact counterexample}

The preceding reductions are sharp.  Take two exceptional profiles of
masses $1/10$ and $9/10$,
\[
 (u_1,v_1)=\left(\frac{624}{625},\frac1{16}\right),
 \qquad
 (u_2,v_2)=\left(\frac1{625},1\right).
\]
Let their first-bulk fibres be complementary indicators of sets of measures
$624/625$ and $1/625$, and let their second-bulk fibres be the constants
$1/16$ and $1$.  Then
\[
 (h_F(i,j))=\begin{pmatrix}624/625&0\\0&1/625\end{pmatrix},
 \qquad
 (h_G(i,j))=\begin{pmatrix}1/256&1/16\\1/16&1\end{pmatrix}.
\]
The density displacements are $609/10000$ and $1/625$.  Substitution into
the exact Atlas $152$ kernel gives the eightfold quadratic matrix
\begin{equation}
 M=\begin{pmatrix}
 34569/20000000&-111/625000\\
 -111/625000&1/78125
 \end{pmatrix},
 \qquad
 \det M=-\frac{2943}{312500000000}<0,
 \label{eq:atlas152-negative-matrix}
\end{equation}
and
\[
 (1/10,9/10)M(1/10,9/10)^{\mathsf T}
 =-\frac{8631}{2000000000}<0.
\]

The independent note
\path{notes/atlas_152_fractional_fibre_local_counterexample.tex} realizes
these fibres as a five-step graphon and evaluates the complete gap at
$\varepsilon=1/3000$.  Its exact density is
$18750093281/37500000000>1/2$, and its exact gap is strictly negative.
Thus Atlas $152$ changes from open to negative; no asymptotic remainder is
used in the accepted certificate.

\section{Hidden-assumption and limitation audit}

The proof uses only finite macro-colouring expansion, Fubini for bounded
functions, the elementary Frechet bounds, and exact polynomial certificates.
It assumes neither regularity nor a finite-step representation of the
graphon.  Shared-neighbour overlaps are retained explicitly and are not
replaced by products of their means.  There is no division, conditioning,
or zero-denominator case.  A common repair is applied only after aggregating
all exceptional profiles.

The positive theorems control the complete quadratic coefficient on the five
stated flat faces and the complete zero-axis mixture cone for four further
rows.  For those four rows they additionally give the exact
arbitrary-kernel endpoint cubic, a uniform theorem on the whole
arbitrary-fibre axis cone, and finally a common radius $1/10000$ for every
measurable two-sided profile.  Thus no scale-dependent profile or kernel
regime remains in this single-exceptional-set model.  This is not a theorem
for several independent exceptional sets or a cut-distance neighbourhood.
The separate two-defect-colouring theorem supplies a uniform $L^\infty$
neighbourhood for the four rows, but it does not subsume the present
order-one small-support statement.  For Atlas $152$, the two-profile indicator,
axis-profile, comparable-mean, equal-displacement, bounded-dispersion, and
comonotone subcones remain valid positive statements, but the complementary
fractional-fibre cone contains the exact counterexample above.  No tensor or
unrelated negative construction is used.

Run
\begin{verbatim}
python codes/needle_higher_order_tests.py
\end{verbatim}
from the repository root for the bounded-memory critical-axis audit used by
the new theorems.  The historical exhaustive suite is available explicitly
as
\begin{verbatim}
python codes/needle_higher_order_tests.py --full
\end{verbatim}
but is not the default because it constructs $12608$ unrelated symbolic
copositivity matrices.  In the full output, the line headed
\texttt{threshold arbitrary-fibre positive pair kernels} reconstructs all
ten endpoint kernels and all sixty simplex certificates using exact rational
arithmetic.  The two-indicator-profile line checks overlap monotonicity and
the six determinant certificates in
Theorem~\ref{thm:atlas152-two-indicator}; the axis-profile line checks
\eqref{eq:atlas152-same-axis} and~\eqref{eq:atlas152-cross-axis}, and the
bounded-dispersion line checks its exact Jensen slack and margin.  The
critical-axis pair line checks \eqref{eq:axis-needle} and all $24$ same-axis,
opposite-axis, and mutual-cell endpoint versions of
\eqref{eq:axis-kernel}; the following rooted-moment line reconstructs the
four raw and repaired identities in
\eqref{eq:axis-endpoint-cubic}, including cancellation of every cross-kernel
moment; the uniform-support line checks
\eqref{eq:axis-endpoint-density-coercivity} and all four exact radius slacks;
the full-axis uniform-support line checks \eqref{eq:full-axis-density},
\eqref{eq:full-axis-coercivity}, the good-set retained factor, and the four
common-radius slacks in Theorem~\ref{thm:full-axis-uniform};
the arbitrary-profile line checks the exact shifted targets, the
one-exceptional-vertex constants, and both case slacks in
Theorem~\ref{thm:full-profile-uniform};
the obstruction line then checks the two exact equal-mean cubic
comparisons.  The
comparable-mean line reconstructs the nine degree-four simplex certificates
and their face-to-face incidence audit; the equal-displacement line
reconstructs its four degree-four certificates and incidence audit.  The
counterexample line reconstructs \eqref{eq:atlas152-negative-matrix}, the
full $5^6$ assignment sum, and the fixed rational witness.  The four
obstruction lines headed
\texttt{Atlas 152} reconstruct the failed $uv$-moment,
total-mass, geometric-mean, and independent-overlap relaxations and
distinguish each from the compatible positive functional.  The comonotone
line verifies the vertex slacks in the pointwise nested-fibre theorem.

\end{document}
